\documentclass[final,a4paper,12pt]{article}

\usepackage{notes}
\setlength\parskip{1ex}
\setlength\parindent{0cm}

\usepackage{float}
\usepackage{mathtools}
\usepackage{amsthm}
\theoremstyle{definition}
\newtheorem{example}{Example}[section]
\newtheorem{problem}{Problem}[section]


\title{Geophysics III}
\author{Complex Numbers Tutorial}
\date{}


\begin{document}
\maketitle

\epigraph{``You would have to be half mad to dream me up.''}
{{\normalfont ---\ Lewis Carroll}, Alice in Wonderland}

\section{Number Sets}
Before we go down the rabbit hole, I need to introduce you to a few bits of
terminology.  Mathematicians especially, but physicists also use these terms and
notation as shorthand to identify groups of numbers with certain
characteristics.  These are identified in Table~\ref{tab:sets}.  The largest
number set is the set of all complex numbers which includes all combinations of
real and imaginary numbers (Figure~\ref{fig:sets}).
\begin{table}[hbt]
    \caption{Number Sets}
    \label{tab:sets}
    \begin{tabularx}{\textwidth}{cX}
        Notation & Number set \\
        $\mathbb{N}$ & Natural (counting) numbers: $1, 2, 3, \ldots$ \\
        $\mathbb{Z}$ & Integers: $\ldots -3, -2, -1, 0 1, 2, 3, \ldots$;
        includes $\mathbb{N}$ \\
        $\mathbb{R}$ & Real numbers: all rational and irrational numbers,
        $(-\infty, \infty)$; rational numbers include all integers and numbers
        that can be expressed as one integer divided by another, while
        irrational numbers are the numbers which cannot be expresses as integer
        ratios (e.g., $\pi$, $e$, $\sqrt{2}$);
        includes $\mathbb{N}$ and $\mathbb{Z}$ \\
        $\mathbb{R}^+$ & Real positive numbers, $(0, \infty)$ \\
        $\mathbb{R}^+_0$ & Real positive numbers and zero, $[0, \infty)$ \\
        $\mathbb{I}$ & Imaginary numbers: all numbers that can be expressed as a
        real number times $\sqrt{-1}$, i.e. numbers that when squared give a
        negative, rather than positive, number \\
        $\mathbb{C}$ & Complex numbers: All numbers containing a real
        ($\mathbb{R}$) and an imaginary part ($\mathbb{I}$)
    \end{tabularx}
\end{table}
\begin{figure}[hbt]
    \centering
    \includegraphics[height=4cm]{figures/number_sets.eps}
    \caption{Relationship between the number sets listed in
    Table~\ref{tab:sets}.  Sets such as $\mathbb{N}$ are entirely captured by
    $\mathbb{Z}$.  All numbers, both real, $\mathbb{R}$ and purely imaginary,
    $\mathbb{I}$, as well as any arbitrary combination of the two, are included
    in the set of complex numbers, $\mathbb{C}$.  The $\mathbb{R}$ and
    $\mathbb{I}$ sets intersect only at $0$ (white point) and therefore are
    orthogonal.}
    \label{fig:sets}
\end{figure}

\section{Introduction}
Mathematicians have problems with things that are undefined.  I can't blame them.
There is something not satisfying when trying to compute a number and instead
you run up against something that is not allowed under the current set of rules.
So what is a mathematician do?

Make new rules and see where they lead.  Take for example the square root
function.  In a more classical scheme, the square root function is
typically defined for positive numbers, i.e.,
\[
    y =
    \begin{cases}
        \sqrt{x} & x \geq 0 \\
        \mathrm{undefined} & \text x < 0 \\
    \end{cases} ,
\]
where $x, y \in \mathbb{R}$.  But what happens when x is negative?  For example,
imagine the case when $y = \sqrt{-1}$.  Mathematicians long ago considered this
question, but initially found no significance beyond mere curiosity; hence, the
name imaginary was given to this group of numbers.  But as it turns out, they
can be used as a compact way to describe waves (we're a bit far from that at the
moment, but we'll get there).

\section{Notation}
We denote the imaginary number by $i$, defined as
\begin{equation}
    i \equiv \sqrt{-1}.
\end{equation}
Hence, $i^2 = -1$.  We can write any negative root as
\begin{align*}
    \sqrt{-b^2} &= \sqrt{(-1)(b^2)}, \\
    &= b \sqrt{-1}, \\
    &= b i,
\end{align*}
where $b \in \mathbb{R}$.  What this tells us is that we can represent an
imaginary number as {\it the} imaginary number, $i$, multiplied by a real number.

We can now extend our definition of the square root to include the root of
negative numbers,
\begin{equation*}
    y = \begin{cases}
        \sqrt{x} & x \geq 0 \\
        i \sqrt{\abs{x}} & x < 0 \\
        \end{cases} ,
\end{equation*}
where $x \in \mathbb{R}$, $\abs{x}$ is the absolute value of $x$, and $y \in
\mathbb{C}$.

Although this tutorial covers complex numbers, dealing with them is not
necessarily complex.  In fact, if you are familiar and comfortable with vectors,
dealing with complex numbers is in some ways very similar.

A complex number has both a real part and a purely imaginary part and can be
written in the form,
\begin{equation}
    \mlbox[ ] {z} {complex number}
    =
    \mlbox[
        labelstyle={xshift=-1cm, below of=m\themathLableNode,below
        of=m\themathLableNode},
        linestyle={out=90,in=270, latex-}]
        {a} {real part}
    +
    \mlbox[
        labelstyle={xshift=1cm, below of=m\themathLableNode,below
        of=m\themathLableNode},
        linestyle={out=90,in=270, latex-}]
        {b i} {imaginary part}
\end{equation}
where $a, b \in \mathbb{R}$.  Just like vector components, the real part and
imaginary part cannot be combined.  There are several other ways to represent a
complex number and we'll get to those in a bit.


\section{The Complex Plane}
Imagine an infinite line (no artists here, it must be straight).  On this line,
we'll add a zero.  It doesn't matter where we place it since the line is
infinite on either side, it just gives us a point of reference from which to
define our positive and negative directions.  If we make everything to the right
positive and everything to the left negative, we can mark out divisions at
integer steps from zero.  This is our real number line and represents all real
numbers.  Between the integers, we have the infinite number of points between
including the rational and irrational numbers.
\begin{figure}[hbt]
    \centering
    \includegraphics[]{figures/real_number_line.eps}
    \caption{The real number line.  Selected numbers from some of the number
    sets are identified.  While example rational and irrational numbers are
    singled out of the real number line, in reality, there are an infinitely
    many numbers between each integer value.}
    \label{fig:real}
\end{figure}

Now imagine a second number line much the same, but each tick mark on this new
line represents an integral step in either the positive, $+i$, or negative,
$-i$, direction.  In between we have all the rational and irrational imaginary
numbers all the same.  We have a $0i$, but $0$ multiplied by any number, $i$
included, is $0$.  Hence, $0$ occurs on both the real and imaginary number lines
indicating that they intersect.  Because real and imaginary numbers cannot be
added, the two number lines must be orthogonal (at right angles to each other).

So now plotting complex numbers is much like plotting (x,y) Cartesian pairs that
you are already familiar with (Figure~\ref{fig:complex}).
\begin{figure}[hbt]
    \centering
    \includegraphics[]{figures/complex_plane.eps}
    \caption{The complex plane, $\mathbb{C}$ includes both the real,
    $\mathbb{R}$, and imaginary number lines, $\mathbb{I}$, which intersect at
    $0 = 0i$.  An example point is shown at $(x,y) = (1,2)$.}
    \label{fig:real}
\end{figure}

\section{Math Operations involving Complex Numbers}
\subsection{Complex Operators}
Note for the imaginary function it returns a real number $b$.  The final
operator that we will define is called the {\it complex conjugate} ($*$).
If $z = x + iy$, then the complex conjugate of $z$ is
\begin{equation}
    z^* = x - iy,
\end{equation}
which simply changes the sign of the imaginary part.  This may seem like a very
silly operator, but, as you will soon see, is vital to computing the magnitude
(norm) of a complex number.

At times, we want to isolate the real part.  To signify that we are only using
the real part we define a function, $\Re$, which returns the real part.  For
example, for a complex number, $z$, written in the form $x + iy$, we want the
function to yield
\begin{equation*}
    \Re(z) = x .
\end{equation*}
Likewise, there we need a function that returns the imaginary part,
\begin{equation*}
    \Im(z) = y .
\end{equation*}
Each of these functions are defined in terms of the complex conjugate,
\begin{align}
    \Re(z) &= \frac{z + z^*}{2} \label{eq:real} \\
    \Im(z) &= \frac{z - z^*}{2i} . \label{eq:imag}
\end{align}

\subsection{Basic Operations}
Many arithmetic operations involving complex numbers are very similar to those
for vectors where the real and imaginary parts are treated as separate
components.  Consider two complex numbers, $z_1 = a + ib$ and $z_2 = c + id$.
To add and subtract $z_1$ and $z_2$,
\begin{align*}
    z_1 + z_2 &= (a + ib) + (c + id) \\
    &= (a + c) + i(b + d) , \\
    z_1 - z_2 &= (a + ib) - (c + id) \\
    &= (a - c) + i(b - d) .
\end{align*}
To multiply, distribute as you might a polynomial and collect terms,
\begin{align*}
    z_1 z_2 &= (a + ib) (c + id) \\
    &= ac + iad + ibc + i^2 bd \\
    &= ac + iad + ibc - bd \\
    &= (ac - bd) + i (ad + bc) . \\
\end{align*}

To divide, we employ the complex conjugate to remove a complex number from the
denominator, e.g.,
\begin{align*}
    \frac{z_1}{z_2} &= \frac{(a + ib)}{(c + id)} \\
    \frac{z_1}{z_2}\frac{z_2^*}{z_2^*} &=
        \frac{(a + ib)}{(c + id)} \frac{(c - id)}{(c - id)} \\
    \frac{z_1 z_2^*}{\abs{z_2}} &= 
        \frac{(ac + bd) + (bc - ad) i)}{(c^2 + d^2)} . \\
\end{align*}

\begin{problem}
Show
\begin{equation*}
    \frac{1}{1 - i} = \frac{1 + i}{2} .
\end{equation*}
\end{problem}

\subsection{Norm}
How big is $4 + 1i$ and is it greater than $2 + 3i$?  We often need to know just
how big a complex number is, but how to tell?  This concept will make a bit more
sense once we define the complex number in a trigonometric representation (e.g.,
Figure~\ref{fig:complex}).  The size (magnitude) of a complex number is called
the norm and is computed as a distance from the origin to the complex point.
For a vector, this is simply done by using the Pythagorean theorem; i.e., if $v
= x \hat{x} + y \hat{y}$,
\begin{align*}
    \abs{v} &= (v \cdot v)^{1/2} \\
    &= [(x \hat{x} + y \hat{y}) \cdot (x \hat{x} + y \hat{y})]^{1/2} \\
    &= (x^2 + y^2)^{1/2} .
\end{align*}
where $\hat{x}$ and $\hat{y}$ are unit vectors in the $x$ and $y$ directions,
respectively.

We wish to define a similar norm for complex numbers, but if you were to say
square each component, we would have the following problem: $x^2 + (iy)^2 = x^2
- y^2$.  Note that this is the difference the squares, not the length we
expected.  For complex numbers, the norm is defined in terms of the number and
its conjugate as
\begin{align*}
    \abs{z} &= (zz^*)^{1/2} \\
    &= [(x + iy)(x - iy)]^{1/2} \\
    &= (x^2 - ixy + ixy  - i^2 y^2)^{1/2} \\
    &= (x^2 - (-1)y^2)^{1/2} \\
    &= (x^2 + y^2)^{1/2} .
\end{align*}


\section{Alternative Representations}
\subsection{Trigonometric Functions}
The $x + iy$ form makes it easy to separate the real and imaginary parts and
plot on the complex plane.  However, as you will see when we get to functions,
the standard form may not be the most straightforward physical description.  To
aid interpretation, we can write a complex number as
\begin{equation}
    z =
    \mlbox[ ] {\abs{z}}{magnitude} 
    (\cos 
    \mlbox[
    labelstyle={xshift=0.85cm, below of=m\themathLableNode,below
        of=m\themathLableNode},
        linestyle={out=90,in=270, latex-}]
    {\phi}{phase angle} + i \sin 
    \mlbox[
    labelstyle={xshift=-0.85cm, below of=m\themathLableNode,below
        of=m\themathLableNode},
        linestyle={out=90,in=270, latex-}]
    {\phi}{}) ,
\end{equation}
where $\phi$ is the phase angle.

This trigonometric representation comes from transformation of a complex number
from $x + iy$ form to a polar coordinate reference frame.  Just as we would for
a Cartesian coordinate pair, we can determine the polar parameters in terms of
the distance from the origin (magnitude) and an angle, which we define as
counterclockwise from the real axis (Figure~\ref{fig:complex}).
\begin{figure}[hbt]
    \centering
    \includegraphics[]{figures/complex_polar.eps}
    \caption{Complex plane and relationships between orthogonal coordinates, a
    circle (polar coordinates), and the natural number.}
    \label{fig:complex}
\end{figure}

To find the magnitude, we compute the norm,
\begin{equation*}
    \abs{z} = (x^2 + y^2)^{1/2} .
\end{equation*}
Note that I use the single bar, $|$, on either side of the variable above to
denote the norm, but it is very common to see a double bar to denote the same
(e.g., $\norm{z}$).  The $\abs{z}$ or $\norm{z}$ are used by some authors to
indicate different things so it is a good idea to determine how the author
defines them.  It may seem as if I have defined the complex number in terms of
itself---or rather it's norm---but in reality, $\abs{z}$ is a real-valued
constant or a real-valued function (i.e., a function with no complex part).

The phase is determined by the ratio of the imaginary to the real parts by
\begin{equation}
    \phi = \tan^{-1} \frac{\Im{z}}{\Re{z}} .
    \label{eq:phase}
\end{equation}

By using Equations~\ref{eq:real} and \ref{eq:imag}, we find the useful
result,
\begin{align*}
    \tan \phi &= \frac{\displaystyle \frac{z - z^*}{2}}
    {\displaystyle \frac{z + z^*}{2}} \\
    \tan \phi &= \frac{z - z^*}{z + z^*} . \\
\end{align*}

To convert from the trigonometric representation to $x + iy$ form,
\begin{align}
    x &= \abs{z} \cos \phi , \\
    y &= \abs{z} \sin \phi .
\end{align}

\begin{problem}
Using $z = x + iy$, show that $\tan \phi = \frac{y}{x}$.
\end{problem}

\begin{problem}
What is the complex conjugate of $z = \abs{z}(\cos \phi + i \sin \phi)$?
\end{problem}



\subsection{Exponential Function}
The exponential form is advantageous because it is perhaps the most compact way
to write a complex number (function) and contains the same easily identifiable
information content as the trigonometric form.  Another reason to write a
complex number in exponential form is that it tends to simplify derivatives.
Recall $\dv{}{x}e^x = e^x$ for $x \in \mathbb{R}$.  For complex exponentials,
$\dv{}{z}e^z = e^z$ for $z \in \mathbb{C}$.  More on this later.

The exponential form is given as
\begin{equation}
    z = \mlbox[]{\abs{z}}{magnitude} e^{i
    \mlbox[labelstyle={xshift=1.2cm},linestyle={out=270,in=90, latex-}]
    {\phi}{phase}} ,
\end{equation}
where $\abs{z}$ is the magnitude and $\phi$ is the phase angle.  The complex
number $z$ need not be constant, but may be any arbitrary function.  Likewise,
the phase may also be a function defined as above in Equation~\ref{eq:phase}.

{\it Proof:}
The exponential form may seem like a long way from our $x + iy$ or the
trigonometric forms so lets explore where this comes from and then we'll
describe a few important values.  We can represent many functions, $f(z)$, about
a given point, $a$, by an infinite series (called a Taylor series) given by
\begin{align*}
    f(z) &= f(a) + \sum_{n = 1}^\infty f^{(n)}(a) \frac{(z - a)^n}{n!} \\
    f(z) &= f(a) + f^{(1)}(a) (z - a) + f^{(2)}(a) \frac{(z - a)^2}{2!}
        + f^{(3)}(a) \frac{(z - a)^3}{3!} + \ldots +
        + f^{(n)}(a) \frac{(z - a)^n}{n!} + \ldots ,
\end{align*}
where $f^{(n)}(a)$ signifies the $n$-th derivative of $f(z)$ evaluated at $a$
and $n!$ is the factorial function for $n$, i.e., $n! = 1 \cdot2 \cdot 3
\cdot\ldots n$.  At $a = 0$ and when $x \in \mathbb{R}$,
\begin{equation*}
    e^{x} = 1 + x + \frac{x^2}{2!} + \frac{x^3}{3!} + \frac{x^4}{4!}
        + \frac{x^5}{5!} + \ldots .
\end{equation*}
Likewise, we can represent the sine an cosine functions as
\begin{align*}
    \cos x &= 1 - \frac{x^2}{2!} + \frac{x^4}{4!} - \ldots \\
    \sin x &= x - \frac{x^3}{3!} + \frac{x^5}{5!} - \ldots .
\end{align*}

Now suppose $x = i \phi$, the exponential function becomes
\begin{equation*}
    e^{i\phi} = 1 + i\phi + \frac{(i\phi)^2}{2!} + \frac{(i\phi)^3}{3!}
        + \frac{(i\phi)^4}{4!} + \frac{(i\phi)^5}{5!} + \ldots .
\end{equation*}
If we employ the following repeating pattern from integer powers of $i$, i.e.,
\begin{align*}
    i^1 &= i \\ 
    i^2 &= -1 \\ 
    i^3 &= -i \\ 
    i^4 &= 1 \\ 
    i^5 &= i \\
    \vdots ,
\end{align*}
we find
\begin{equation*}
    e^{i\phi} = 1 + i\phi - \frac{\phi^2}{2!} - \frac{i \phi^3}{3!}
        + \frac{\phi^4}{4!} + \frac{i \phi^5}{5!} - \ldots .
\end{equation*}
Using the same relationships above for sine and cosine ($x = \phi$), we find
\begin{equation*}
    \cos \phi + i \sin \phi = 
        1 + i\phi - \frac{x^2}{2!} - \frac{i \phi^3}{3!} + \frac{x^4}{4!} 
        + \frac{i \phi^5}{5!} - \ldots .
\end{equation*}
And thus it is shown that
\begin{equation}
    e^{i \phi} = \cos \phi + i \sin \phi .
\end{equation}
Q.E.D.

\begin{problem}
Using the trigonometric representation, show
\begin{align*}
    e^{i0} & = 1 \\
    e^{i \frac{\pi}{4}} & = \frac{\sqrt{2}}{2} + \frac{\sqrt{2}}{2} i \\
    e^{i \frac{\pi}{2}} & = i \\
    e^{i \pi} & = -1 \\
    e^{i \frac{3\pi}{2}} & = -i \\
    e^{i n \pi} &=
    \begin{cases}
        1 & n\ \text{is even} \\
        -1 & n\ \text{is odd} \\
    \end{cases}
\end{align*}
\end{problem}

\begin{problem}
What is the complex conjugate of $e^{i \phi}$?
\end{problem}


\section{Complex Functions}
Some of the functions that we will examine in class are complex functions
(specifically waves).  A function similar to one we will discuss in class can be
written
\begin{equation}
    E_x (z) = E_0 \exp \left[-\frac{(1 + i)}{\sqrt{2}} z \right] .
    \label{eq:cfunc}
\end{equation}
In this example, $z$ is a spatial coordinate and no longer a complex number.
The function is shown in Figure~\ref{fig:cfunc3}.  The resultant function is a
spiral in the complex plane that exponentially approaches zero as $z$ increases.
Functions in 3-dimensional plots can be relatively difficult to interpret so we
often represent them parametrically\footnote{Parametric equations are ones in
which the coordinates are represented as functions of a single variable.
Typically, these are represented in polar notation.  For example, for a circle
with radius $r$, $r^2 = x^2 + y^2$, can be represented parametrically as a
function of a single parameter, $\theta$,
\begin{align*}  
    x &= r \cos \theta \\
    y &= r \sin \theta .
\end{align*}}.
In the case of complex functions, the parametric
function is plotted on the complex plane.
\begin{figure}[hbt]
    \centering
    \includegraphics[]{figures/complex_func3.eps}
    \caption{A complex function (Equation~\ref{eq:cfunc}) of $z$.  $E_0$ is
    assumed to be $1$.   The real and imaginary parts as a function of $z$ are
    shown in red and as a projection onto the complex plane in blue.
    Figure~\ref{fig:cfunc} shows the projection onto the complex plane.}
    \label{fig:cfunc3}
\end{figure}

We've already discussed the complex plane.  Plotting a complex function is no
different than plotting an ordinary set of $x$ and $y$ components, in this case
real and imaginary.  The difference being the representation of the function not
as a typical input versus output but as a parametric function, that is
\begin{align*}
    x &= \Re(E_x(z)) \\
    y &= \Im(E_x(z)) ,
\end{align*}
where $x$ are the real values of the function and $y$ are the imaginary values
of the function.  One approach is to rewrite our equation in $x + iy$ form so
that the real and imaginary parts can be easily separated, which we can do via
the trigonometric representation.

To write the function in Equation~\ref{eq:cfunc} in $x + iy$ form,
\begin{align*}
    E_x (z) &= E_0 \exp \left[ -\frac{(1 + i)}{\sqrt{2}} z \right] \\
    &= E_0 \exp \left( -\frac{1}{\sqrt{2}} z \right) 
        \exp \left( -\frac{i}{\sqrt{2}} z \right) \\
    &= E_0 \exp \left( -\frac{1}{\sqrt{2}} z \right) 
        \left[ \cos \left( -\frac{1}{\sqrt{2}} z \right) + 
        i \sin \left( -\frac{1}{\sqrt{2}} z \right) \right] .
\end{align*}
We can now define our parametric equations:
\begin{align}
    \Re(E_x(z)) &= E_0 \exp \left( -\frac{1}{\sqrt{2}} z \right) 
        \cos \left( -\frac{1}{\sqrt{2}} z \right) ,\\
    \Im(E_x(z)) &= E_0 \exp \left( -\frac{1}{\sqrt{2}} z \right) 
        \sin \left( -\frac{1}{\sqrt{2}} z \right) .
\end{align}
The parametric functions produce the projection of the spiral on the complex
plane (Figure~\ref{fig:cfunc}).
\begin{figure}[hbt]
    \centering
    \includegraphics[]{figures/complex_func.eps}
    \caption{The parameterized plot of the complex function in
    Equation~\ref{eq:cfunc}.  The spiral nature of this equation is clearly
    evident in this plot.  We can describe each point along the plot in terms of
    two parameters, a distance from the origin (norm) and an angle
    counterclockwise from the real axis (phase).}
    \label{fig:cfunc}
\end{figure}

A complex function can sometimes be a bit difficult to interpret---we live in a
real space, not an imaginary one---so it can be advantageous to write our
equation as real numbers that we can ascribe physical meaning to our function,
namely a magnitude and a phase.  We have already found these parameters
inadvertently on our way to a trigonometric representation,
\begin{align}
    \abs{E_x (z)} &= E_0 \exp \left[ -\frac{z}{\sqrt{2}} \right] ,\\
    \phi &= -\frac{z}{\sqrt{2}} . \\
\end{align}
The results are shown in Figure~\ref{fig:cfunc2}.  Note $2\pi$ has been added to
the phase so that it is positive.
\begin{figure}[hbt]
    \centering
    \includegraphics[]{figures/complex_func2.eps}
    \caption{Representation of the complex function in Equation~\ref{eq:cfunc}
    as a magnitude (a) and phase (b) as a function of the independent variable
    $z$.}
    \label{fig:cfunc2}
\end{figure}
For the parametric equations, the magnitude is given by
\begin{align*}
    \abs{E_x(z)} &= (\Re(E_x)^2 + \Im(E_x)^2)^{1/2} \\
    &= \left[
        E_0^2 \exp \left( -\frac{2}{\sqrt{2}} z \right) 
            \cos^2 \left( -\frac{1}{\sqrt{2}} z \right)
        + E_0^2 \exp \left( -\frac{2}{\sqrt{2}} z \right) 
            \sin^2 \left( -\frac{1}{\sqrt{2}} z \right)
    \right]^{1/2} \\
    &= E_0 \exp \left( -\frac{1}{\sqrt{2}} z \right) 
        \left[ \cos^2 \left( -\frac{1}{\sqrt{2}} z \right)
            + \sin^2 \left( -\frac{1}{\sqrt{2}} z \right)
        \right]^{1/2} \\
    &= E_0 \exp \left( -\frac{1}{\sqrt{2}} z \right) 
        ( 1 )^{1/2} \\
    &= E_0 \exp \left( -\frac{1}{\sqrt{2}} z \right) , \\
\end{align*}
and the phase is given by
\begin{align*}
    \phi &= \tan^{-1} \left[\frac{\Im(E_x)}{\Re(E_x)} \right] \\
    &= \tan^{-1} \left[\frac{E_0 \exp \left( -\frac{1}{\sqrt{2}} z \right) 
        \sin \left( -\frac{1}{\sqrt{2}} z \right)}
        {E_0 \exp \left( -\frac{1}{\sqrt{2}} z \right) 
        \cos \left( -\frac{1}{\sqrt{2}} z \right)} \right] \\
    &= \tan^{-1} \left[\frac{\sin \left( -\frac{1}{\sqrt{2}} z \right)}
        { \cos \left( -\frac{1}{\sqrt{2}} z \right)} \right] \\
    &= \tan^{-1} \left[\tan \left( -\frac{1}{\sqrt{2}} z \right) \right] \\
    &=  -\frac{1}{\sqrt{2}} z .
\end{align*}


\section{Derivatives of Complex Functions}
I won't have you perform any derivatives of complex functions as part of the
course.  But for those that are interested, I've included this brief
(non-comprehensive) section that will give you a basis in the idea of complex
derivatives.  Taking derivatives of complex functions is not necessarily much
more complex than derivatives of real-valued functions.  The sum, product, chain
rule, etc. all still hold.  However, there are extra conditions that must be met
in order for the derivative to exist.

Consider a complex function, $f(z) = u(x,y) + i v(x,y)$, which we construct
using two functions $u$ and $v$.  Recall from our calculus discussions that the
derivative at a point, $z_0$, is defined as
\begin{equation*}
    \left. \dv{}{z} f(z) \right|_{z = z_0} =
        \lim_{\Delta z \to 0} \frac{f(z_0 + \Delta z) - f(z_0)}{\Delta z}
\end{equation*}
Since $\Delta z = \Delta x + i \Delta y$, $\Delta z \to 0$ can occur by either
$\Delta x \to 0$ or $i \Delta y \to 0$.  Hence the derivative of a complex
function, $f(z)$, may be performed by either
\begin{align}
    \dv{}{z} f(z) &= \pdv{u}{x} + i \pdv{v}{x} ,\label{eq:dzdx}\\
    \dv{}{z} f(z) &= -i\pdv{u}{y} + \pdv{v}{y} .\\
    \label{eq:dzdy}
\end{align}
These equations leads to the Cauchy-Riemann equations,
\begin{gather}
    \pdv{u}{x} = \pdv{v}{y} ,\\
    \pdv{u}{y} = -\pdv{v}{x} ,
\end{gather}
implying that we have somewhat of a choice in how we determine the derivative of
our complex function.  The Cauchy-Riemann equations also show that the functions
$u$ and $v$ must be related in this manner at every point $z_0$, and if not, the
function is not differentiable.

Let's take a look at $e^z$ and show that it's derivative is in fact itself.
Using the standard from for a complex number,
\begin{align*}
    e^z &= e^{x + i y} \\
    &= e^x e^{i y} \\
    &= e^x (\cos y + i \sin y) \\
    &= e^x \cos y + e^x i \sin y .
\end{align*}
Since we can choose $x$ or $y$ to take our derivative, we'll take $x$ for
simplicity first,
\begin{align*}
    \dv{}{z} e^z &= \pdv{}{x}(e^x \cos y) + i \pdv{}{x} (e^x \sin y) \\
    &= e^x \cos y + i e^x \sin y \\
    &= e^z .
\end{align*}
But we aren't done.  We must show that the Cauchy-Riemann equations hold or that
both relationships in Equations~\ref{eq:dzdx},
\begin{align*}
    \pdv{}{x} (e^x \cos y) &= \pdv{}{y} [e^x \sin y] \\
    e^x \cos y &= e^x \cos y ,
\end{align*}
and Equation \ref{eq:dzdy} yield the same
result,
\begin{align*}
    \pdv{}{y} (e^x \cos y) &= -\pdv{}{x} [e^x \sin y] \\
    - e^x \sin y &= - e^x \sin y .
\end{align*}
Since the Cauchy-Riemann conditions are met and $e^z$ is continuous, the
derivative of $e^z$ is in fact itself.

\end{document}
